\begin{abstract}
% It is well established that language models obey scaling laws, i.e. there exists power-law dependence between loss and the number of parameters and training steps. We take scaling laws one step further by introducing the granularity parameter, which gives control over the number of experts and size of a single expert in a way that the number of FLOPs used by the model is constant (a more granular model consists of smaller experts). We show that increasing granularity helps to achieve better results and also we provide a formula for compute optimal granularity, taking into account routing overhead existing in all Mixture of Experts models. For several compute budgets for training, we train granularity-optimal models and compare them against baselines (dense transformer, and routing models), showing a clear improvement.

Mixture of Experts (MoE) models have emerged as a primary solution for reducing the computational cost of Large Language Models. In this work, we analyze their scaling properties, incorporating an expanded range of variables.  Specifically, we introduce a new hyperparameter, granularity, whose adjustment enables precise control over the size of the experts. Building on this, 
we establish scaling laws for fine-grained MoE, taking into account the number of training tokens, model size, and granularity.  Leveraging these laws, we derive the optimal training configuration for a given computational budget. Our findings not only show that MoE models consistently outperform dense Transformers but also highlight that the efficiency gap between dense and MoE models widens as we scale up the model size and training budget. Furthermore, we demonstrate that the common practice of setting the size of experts in MoE to mirror the \ff{} layer is not optimal at almost any computational budget. 
\noindent\let\thefootnote\relax\footnote{Contributions: Jakub implemented fine-grained MoE, ran experiments, and oversaw the course of the project. Jan designed and implemented the scaling laws, also optimized and tuned the fine-grained MoE implementation.
Kamil A. provided significant advice on many aspects of the project. Maciej experimented with the block design and, with Michał, provided considerable technical support. Szymon, Kamil C., Krystian, and Tomasz contributed to the project and the engineering in various ways. Marek, along with Piotr, provided high-level scientific advice. Sebastian came up with the initial idea, started the project, and supervised it while setting the research direction and leading experiments and analyses. Correspondence to \textless{}s.jaszczur@uw.edu.pl\textgreater{}. $^{\ast}$~Equal contribution. 
}

%albo: Consequently, we show that MoE models always outperform dense transformers. Furthermore, the efficiency gap between dense and MoE models widens as we scale up the model size and training budget.

%Mixture of Experts (MoE) models have emerged as a primary solution for reducing the computational cost of Large Language Models. In this work, we analyze their scaling properties, highlighting certain arbitrary assumptions present in the existing literature. In particular, we introduce a new hyperparameter, granularity, the modification of which allows for the optimal adjustment of the size of experts. Subsequently, we present scaling laws for fine-grained MoE, taking into account the number of training tokens, model size, and granularity. Using these scaling laws, we derive the optimal training configuration for a given computational budget. Furthermore, in contrast with previous works, we demonstrate that the gap in efficiency between dense and MoE models grows as we scale up the model size and training budget.\jk{dodać TO}


% \sj{Can we base it on some standard format, e.g. ICLR? For example, the margins here are extraordinarily small; it looks weird.} \jk{spoko, możesz podmienić, tutaj jest szablon jak w MoE-Mambie}
% In recent years Mixture of Experts (MoE) has emerged as a primary remedy for  reducing computational cost while scaling Large Language Models. Till this paper  MoE architectures assumed that each feed forward layer in a transformer network is
% replaced with $\nexpert$ experts with the size of the original layer. Here, we observe  that this implicit assumption resulted in sub-optimal architectures, and thus calls for further thoughtful study that will validate the usefulness of MoE for scaling 
% LLMs further. This work takes on this goal delivering the following two novel results:
% \begin{itemize}
% \item we introduce the notion of {\em granularity}, i.e., smaller experts, extending  the space of hyperparameters for the model architecture and allowing to seek for optimal granularity,
% \item we derive empirical scaling laws demonstrating that minimizing perplexity for given computational budget requires higher granularity for larger models.
% \end{itemize} 

%Mixture of Experts (MoE) has become a standard technique for scaling Large Language Models without much increase in the number of floating point operations required for their inference. However, until now, an implicit assumption about the size of each expert ($\dexpert = 4\dmodel$) 
%\ka{what assumption?} 
%in MoE has resulted in sub-optimal architectures. In this work, we introduce the notion of granularity, which generalizes current MoE approaches and derive empirical scaling laws, that allow choosing optimal granularity, which achieves lower perplexity than current approaches for a given computational budget. \ka{a little more here}
% \mk{maybe mention that this models achieves lower complexity}. 
% We also introduce a modification to an established MoE architecture (Expert Choice), making it more versatile and show that it can be fine-tuned to work in the Token Choice regime. 
\end{abstract}